% Part: incompleteness
% Chapter: introduction
% Section: historical-background

\documentclass[../../../include/open-logic-section]{subfiles}

\begin{document}

\olfileid{inc}{int}{bgr}

\olsection{تاریخی پس منظر}

ایس حصے وچ اسی مختصر طور اُتے اوہناں تاریخی پیش رفتاں دی گل
کراں گے جیہڑیاں نامکملیت دے قضیاں دا پس منظر سمجھّن وچ مدد
دیندیاں نیں۔ خاص طور اُتے، پہلے ریاضیاتی منطق دی تاریخ دا
سرسری خاکہ دیواں گے۔ فیر ریاضی دیاں بنیاداں دی تاریخ
بارے کجھ گلّاں کراں گے۔

\begin{digress}
  «ریاضیاتی منطق» دی ترکیب دے دو مطلب ہو سکدے نیں۔ «ریاضیاتی»
  نوں موضوع دی صفت سمجھو تاں مراد «ریاضی دی منطق» اے، یعنی
  ریاضی وچ استدلال دے اصول۔ جے ایہنوں طریقے دی صفت سمجھو
  تاں مراد «منطق دی ریاضی» اے، یعنی استدلال دے اصولاں دا
  ریاضیاتی مطالعہ۔ اگلے بیان وچ ریاضیاتی منطق دے ایہ دونیں
  مطلب آؤندے نیں، کئی وار اکّو ویلے۔
\end{digress}

منطق دا باقاعدہ مطالعہ ارسطو توں شروع ہویا منیا جاندا اے۔
اوہ لگ بھگ 384--322 \textsc{قبلِ مسیح} وچ جیویا۔ اوہ دیاں
\emph{زمرے (Categories)}، \emph{پہلے تجزیے (Prior
  analytics)} تے \emph{بعد دے تجزیے (Posterior analytics)}
علمی استدلال دے اصولاں دا باقاعدہ مطالعہ کردیاں نیں۔
ایہناں وچ قیاس، یعنی مقرر مقدمیاں توں نتیجہ کڈھن والی
صورت، دا وی پورا تے منظم مطالعہ اے۔

ارسطو دی منطق قرونِ وسطیٰ وچ مدرسی فلسفے اُتے حاوی رہی۔
ایتھوں تک کہ اٹھارہویں صدی وچ وی کانٹ دا خیال سی کہ
ارسطو دی منطق مکمل اے تے ایہنوں بدلّن دی لوڑ نہیں۔
پر قیاس دا نظریہ ریاضیاتی استدلال دے سطحی ترین پہلو توں
اگے کوئی ٹھیک نمونہ نہیں بنا سکدا۔ اوہدے توں اک صدی
پہلاں، نیوٹن دے ہم عصر لائبنز نے منطقی استدلال لئی
اک مکمل «حساب» دا تصور کیتا سی۔ اوہنے ایداں دا حساب
تیار کرن ول کجھ ابتدائی قدم وی چُکے، تے بنیادی طور اُتے
قضیاتی منطق دی اک صورت بیان کیتی۔

انیویں صدی منطق لئی فیصلہ کن موڑ سی۔ 1854 وچ
جارج بُول نے \emph{سوچ دے قانوناں (The Laws of Thought)}
لکھی، جیہدے وچ قضیاتی منطق دا مفصل الجبری مطالعہ اے۔
ایہ مطالعہ اج دی پیش کش توں بہتا وکھرا نہیں۔ 1879 وچ
گوٹلوب فریگے نے اپنی \emph{Begriffsschrift}
(«تصوری تحریر») شائع کیتی۔ اوہنے قضیاتی منطق وچ
مسوّراں تے تعلق جوڑے۔ ایس طرح ایہدے وچ
پہلے درجے دی منطق وی شامل سی۔ حقیقت وچ فریگے دے
منطقی نظاماں وچ اعلیٰ درجے دی منطق تے ہور وی کجھ
شامل سی۔ اپنی \emph{حساب دے بنیادی قانوناں
  (Basic Laws of Arithmetic)} وچ فریگے نے ایہ
دکھاون دی کوشش کیتی کہ سارا حساب اوہدی
Begriffsschrift وچ صرف منطقی مفروضیاں توں اخذ
ہو سکدا اے۔ بدقسمتی نال اوہ مفروضے متضاد نکلے،
جیویں رسل نے 1902 وچ دکھایا۔ پر متضاد مسلمے
نوں پاسے رکھ دیئیے تاں فریگے نے جدید منطق لگ بھگ
اکیلے ہی ایجاد کر دتی؛ ایہ حیران کن کارنامہ سی۔
بُول توں بعد الجبری سوچ رکھنے والے ہور مفکراں،
جنہاں وچ پئرس تے شروڈر شامل سن، نے وی
مسوّری منطق آزاد طور اُتے تیار کیتی۔

ہن ریاضی دیاں بنیاداں وچ پیش رفت ول آئیے۔
ظاہر اے کہ منطق ریاضی وچ اہم کم کردی اے، ایس لئی
ایہ پیش رفت منطق دیاں اوہناں تبدیلیاں نال
جیہڑیاں اسیں ہُن تک بیان کیتیاں نیں،
کافی جڑی ہوئی اے۔ مثال لئی فریگے نے اپنی
منطق ایس واضح مقصد نال بنائی سی کہ ساری
ریاضی صرف اوس دے منطقی ڈھانچے اُتے کھڑی
ہو سکدی اے۔ خاص کر اوہ دکھانا چاہندا سی
کہ ریاضی تجربے توں آزاد \emph{تجزیاتی}
سچائیاں اُتے مشتمل اے، نہ کہ تجربے توں
آزاد \emph{ترکیبی} سچائیاں اُتے، جیویں
کانٹ سمجھدا سی۔

بہت سارے لوک ریاضی دے باقاعدہ جنم نوں یونانیاں نال
جوڑدے نیں۔ لگ بھگ 300 قبلِ مسیح لکھی اقلیدس دی
\emph{عناصر (Elements)} پہلے ہی یونانی ریاضی دی
اک پختہ مثال اے، جیہدے وچ دلیل دی پختگی تے ٹھیک پن اُتے
زور اے۔ اقلیدس دی \emph{عناصر} وچ تعریفاں تے
ثبوت اج وی ہائی سکولاں لئی ہندسے دیاں درسی کتاباں وچ لگ بھگ
اوہو جہے موجود نیں (جتھے ہائی سکولاں وچ ہندسہ ہن وی
پڑھایا جاندا اے)۔ ریاضیاتی استدلال دے ایس نمونے
نوں صرف ریاضی وچ نہیں بلکہ فلسفے دیاں شاخاں وچ وی
سخت استدلال دی مثال سمجھیا گیا اے۔ (اسپینوزا نے
تاں اخلاقی تے مذہبی دلیلیں وی اقلیدسی طرز وچ پیش
کیتیاں؛ اوہ ویکھ کے حیرانی ہوندی اے!)

سترہویں صدی وچ نیوٹن تے لائبنز نے حسابان ایجاد کیتا۔
(ایہ جھگڑا صدیاں تک چلدا رہیا کہ پہلاں ایہ کم کس نے کیتا۔ اج
اکثر محقق سمجھدے نیں کہ دوناں دی پیش رفت وڈی حد
تک اک دوجے توں آزاد سی۔) حسابان وچ، مثال لئی،
لامتناہی حد تک چھوٹیاں مقداراں دے لامتناہی جوڑاں بارے
استدلال ہوندا اے۔ ایہناں خاصیتاں دی وجہ توں
بشپ برکلے نے تنقید کیتی: اوہدے مطابق رب اُتے
یقین اوہدے ویلے دی ریاضی توں گھٹ معقول نہیں سی۔
اٹھارہویں صدی وچ حسابان دے طریقے وڈے پیمانے
اُتے ورتے گئے؛ مثال لئی لیونارڈ اویلر نے
لامتناہی جوڑاں والے حساب نال حیران کن
نتیجے کڈھے۔

انیویں صدی وچ ریاضی داناں نے حسابان نوں ہور مضبوط
بنیاد اُتے رکھ کے برکلے دی تنقید دا جواب دین دی
کوشش کیتی۔ کوشی، وائرشٹراس، بولزانو تے ہوراں دی
کوششاں توں حد، استمرار، مشتق گیری تے تکامل دیاں
اج والیاں تعریفاں «ایپسلون تے ڈیلٹا» دی مدد نال
بنیاں؛ یعنی اوہناں وچ لامتناہی حد تک چھوٹی
مقداراں دا کوئی حوالہ نہ رہا۔ ایس صدی دے پچھلے
حصے وچ ریاضی داناں نے ہور اگے ودھ کے حسابان
دے سارے پہلو، حقیقی اعداد سمیت، قدرتی عدداں
دی مدد نال سمجھاون دی کوشش کیتی۔ (کرونیکر:
«رب نے پورے عدد بنائے؛ باقی سب انسان دا کم اے۔»)
1872 وچ ڈیڈیکنڈ نے «استمرار تے غیر ناطق اعداد»
(Continuity and the irrational numbers) لکھی،
جیہدے وچ اوہنے ناطق عدداں دے سیٹاں راہیں
حقیقی اعداد دی «تعمیر» دکھائی (تے تُسی جان دے
او کہ ناطق عدداں نوں قدرتی عدداں دے جوڑیاں
وجوں سمجھیا جا سکدا اے)۔ 1888 وچ اوہنے
«Was sind und was sollen die Zahlen» (موٹے
طور اُتے، «قدرتی عدد کیہ نیں تے اوہناں نوں
کیہ ہونا چاہیدا اے؟») لکھی، جیہدا مقصد
قدرتی عدداں نوں خالص «منطقی» لفظاں وچ
سمجھانا سی۔ 1887 وچ کرونیکر نے
«\"Uber den Zahlbegriff» («عدد دے تصور بارے»)
لکھی، جیہدے وچ اوہنے ساریاں ریاضیاتی
اشیاں نوں صحیح عدداں دے ذریعے بیان کرن
دی گل کیتی۔ 1889 وچ جوزیپے پیانو نے
قدرتی عدداں لئی رسمی، علامتی مسلّمے دتے۔

انیویں صدی دے آخر وچ لامتناہی نال برتاؤ
وچ وی اک نویں دلیری آئی۔ اوہدے توں پہلاں
لامتناہی اشیاں تے ساختاں، مثلاً قدرتی عدداں
دا سیٹ، احتیاط نال ورتیاں جاندیاں سن۔
«لامتناہی تعداد وچ» دا مطلب «جنے چاہو اونے»
سمجھیا جاندا سی، تے «حد ول جانا» دا مطلب
«جنا چاہو اونا نیڑے آنا»۔ پر جارج کانٹور
نے دکھایا کہ لامتناہی نوں اپنے پورے معنی
وچ وی لیا جا سکدا اے۔ کانٹور، ڈیڈیکنڈ
تے ہوراں دے کم نے ریاضی نوں سیٹاں دے
عام نظریے راہیں سمجھّن دا اوہ طریقہ متعارف
کراؤن وچ مدد دتی جیہڑا اج وڈے پیمانے
اُتے قبول اے۔

ایہ گلاں سانوں ویہویں صدی دی منطق تے
ریاضی دیاں بنیاداں دی پیش رفت تک لے
آندیاں نیں۔
% PNB-CMP-045: the dated 1902 event is Russell's letter to Frege, not the date of discovery.
1902 وچ رسل نے فریگے نوں اوس دے منطقی
نظام وچ ملے تضاد بارے خط لکھ کے آگاہ
کیتا۔ 1904 وچ زرمیلو نے «انتخاب دے
مسلّمے» دی مدد نال کانٹور دا حُسنِ ترتیب
دا اصول ثابت کیتا۔ ایس مسلّمے دے جائز
ہون بارے کافی بحث ہوئی۔ 1910 تے 1913
دے وچکار رسل تے وائٹ ہیڈ دی
\emph{Principia Mathematica} دیاں تین جلدّاں
شائع ہوئیاں، جنہاں نے ریاضی نوں منطقی
بنیاداں اُتے کھڑا کرن دا فریگے والا
منصوبہ اگے ودھایا۔ پر رسل تے وائٹ ہیڈ
نوں دو ایہو جہے اصول مننّے پئے جنہاں نوں
خالص منطقی کہہ کے جائز ٹھہرانا اوکھا سی:
لامتناہیت دا اک مسلمہ تے «قابلِ تقلیل
ہون» دا اک مسلمہ۔ 1900 دی پہلی دہائی
وچ پوانکارے نے ریاضی وچ «پورے مجموعے
اُتے انحصار کرن والیاں تعریفاں»
(impredicative definitions)
دی ورتوں اُتے تنقید کیتی۔ 1910 دی
دہائی وچ براوَر نے ساری ریاضی نوں
«وجدانیت پسند» بنیاد اُتے دوبارہ
کھڑا کرن دی تجویز پیش کرنی شروع
کیتی، جیہڑی تیسرے امکان دے اخراج
دا قاعدہ~($!A \lor \lnot !A$)
نہیں ورتدی۔

سچ مچ عجیب زمانہ سی!{} ساری ریاضی
نوں منطق وچ گھٹا دین والے منصوبے
نوں ہن «منطقیت» آکھیا جاندا اے؛
اُپر بیان کیتیاں اوکھاں کارن ایہنوں
عام طور اُتے ناکام سمجھیا جاندا اے۔
ریاضی نوں وجدانیت پسند ذہنی تعمیراں
راہیں بنانے والے منصوبے نوں «وجدانیت»
آکھدے نیں، تے عام ریاضی اُتے ایہدیاں
پابندیاں حد توں ودھ سخت سمجھیاں
جاندیاں نیں۔ صدی دے موڑ اُتے ڈیوڈ
ہلبرٹ، ہر ویلے دے سب توں اثر
رکھن والے ریاضی داناں وچوں اک،
کانٹور تے ڈیڈیکنڈ دے نویں تجریدی
طریقیاں دا پکا حامی سی: «کانٹور نے
ساڈے لئی جیہڑی جنت بنائی اے،
کوئی سانوں اوتھوں نہیں کڈھ سکے گا۔»
نال ہی اوہ ایہناں نویں طریقیاں
(جنہاں نوں ہن عجیب طور اُتے
«کلاسیکی» آکھیا جاندا اے) اُتے
بنیادی تنقید بارے حساس سی۔ اوہنے
دونیں فائدے اکٹھے رکھن دا اک راہ
سجھایا:
\begin{enumerate}
\item کلاسیکی طریقیاں نوں رسمی
  مسلّماں تے قاعدیاں نال بیان کرو؛
  ریاضیاتی سوالاں نوں مسلّمی
  نظام وچ !!{formula}s بنا کے پیش کرو۔
\item بھروسے جوگ «متناہی» طریقیاں نال
  ثابت کرو کہ ایہ رسمی اخذی نظام
  بے تضاد نیں۔
\end{enumerate}

پہلا مقصد حاصل کرن ول ہلبرٹ دا
کم بہت اگے ودھیا۔ 1899 وچ اوہنے
اپنی مشہور کتاب
\emph{ہندسہ دیاں بنیاداں (Foundations of geometry)}
وچ ہندسہ لئی ایہ کم کیتا سی۔ بعد
دے سالاں وچ اوہ تے اوس دے کئی
شاگرداں تے ساتھیّاں نے ریاضی
دے ہور میداناں وچ وی ہندسہ
والا کم کرن دی کوشش کیتی۔
ہلبرٹ نے خود حساب تے تجزیہ
لئی مسلّماتی نظام دتے۔ زرمیلو
نے سیٹ نظریے دے مسلّمے بنائے،
جنہاں نوں بعد وچ فرینکل،
سکولم، فان نوئمان تے ہوراں
نے اگے ودھایا۔ 1920 دی
دہائی دے وچکار «ساری» ریاضی
نوں مسلّماتی شکل دین دا دعویٰ
کرن والے دو طریقے موجود سن:
رسل تے وائٹ ہیڈ دی
\emph{Principia mathematica}،
تے اوہ سیٹ نظریہ جیہڑا
زرمیلو--فرینکل سیٹ نظریے
دے ناں نال مشہور ہویا۔

1921 وچ ہلبرٹ نے ایہناں
نظاماں دی بے تضادی ثابت
کرن دا ہدف رکھ کے تحقیقی
منصوبہ شروع کیتا۔ ایس
کم وچ اوس دے کئی شاگرداں،
خاص کر برنائز، ایکرمان
تے بعد وچ گینٹزن، نے
اوہدی مدد کیتی۔ بنیادی
خیال ایہ سی کہ ریاضی وچ
تضاد دا !!{derivation} ہو سکّن
دا سوال علامتاں دیاں ممکن
لڑیاں بارے اک امتزاجی مسئلے
دی صورت وچ رکھیا جاوے۔ یعنی
جملیاں دیاں اوہ لڑیاں ویکھیاں
جان جیہڑیاں حساب،
تجزیے یا سیٹ نظریے دے
مسلّماتی نظام توں،
مثلاً، $!A \land
\lnot !A$ دے صحیح
!!{derivation} ہون دیاں
شرطاں پوریاں کردیاں
ہون۔ ایسی علامتی لڑی
ناممکن ہون دا ثبوت
خود اک ریاضیاتی ثبوت
ہووے گا، ایس لئی
ایہناں مسلّماتی نظاماں
وچ رسمی طور اُتے
بیان کیتا جا سکدا
ہووے گا۔ دوجے لفظاں
وچ، کوئی جملہ~$\OCon$
ہووے گا جیہڑا کہے
گا کہ، مثال لئی،
حساب بے تضاد اے۔
ہور ایہ کہ متعلقہ
نظاماں وچ ایہ جملہ
ثابت وی ہو سکنا
چاہیدا اے، خاص کر
جے اوس دا ثبوت
بہت محدود «متناہی»
وسیلے ہی منگدا ہووے۔

دوجا ہدف، یعنی بنائے
ہوئے مسلّماتی نظاماں
نال ہر ریاضیاتی سوال
دا فیصلہ ہونا، دو طرحاں
نال ٹھیک ٹھیک بیان
ہو سکدا اے۔ پہلا بیان
ایہ اے: ریاضی دے مسلّماتی
نظام دی زبان دے ہر
جملے~$!A$ لئی یا
تاں~$!A$ یا~$\lnot !A$
اوس دے مسلّماں توں
ثابت ہو سکدا ہووے۔
جے ایہ درست ہووے
تاں کوئی جملہ ایہو
جیہا نہ ہووے گا
جیہنوں انہاں مسلّماں
توں نہ ثابت کیتا
جا سکے نہ رد؛
ایہ مسلّمے ہر سوال
دا فیصلہ کردے
ہون گے۔ ایس خاصیت
والے مسلّماتی نظام
نوں \emph{کامل}
آکھدے نیں۔ بے شک
کسے دتے جملے لئی
ایہ طے کرنا ہن وی
اوکھا ہو سکدا اے
کہ دوناں وچوں
کیہڑا متبادل ثابت ہو سکدا اے۔
پر اصولی طور اُتے
ایہ معلوم کرن دا
کوئی طریقہ ہونا
چاہیدا اے۔ ہلبرٹ
دے زیرِ غور مسلّماتی
تے !!{derivation}
نظاماں لئی حقیقت
وچ کاملیت توں
ایہو جیہا طریقہ
موجود ہونا نکلدا
سی، بھاویں ہلبرٹ
نوں ایس دا احساس
نہیں سی۔ سوال
سمجھّن دا دوجا
طریقہ ایس توں
ودھ زورآور شرط
اے: کوئی مشینی،
حسابی طریقہ ہووے
جیہڑا دتے جملے~$!A$
لئی طے کرے کہ
اوہ مسلّماں توں
اخذ ہون جوگ اے
یا نہیں۔

% PNB-CMP-046: restrict Gödel's claims to effective arithmetic-rich consistent theories; own consistency is unprovable, not automatically independent.
1931 وچ گوڈل نے
«نامکملیت دے دو قضیے»
ثابت کیتے، جنہاں
نے دکھایا کہ ایہ
منصوبہ اپنی مطلوبہ
صورت وچ کامیاب
نہیں ہو سکدا۔
ریاضی دا ہر مسلّماتی
نظام نہیں بلکہ
ہر اوہ بے تضاد،
مؤثر طور اُتے
مسلّماں نال پیش
کیتا جا سکّن
والا تے حساب
لئی کافی طاقتور
نظام نامکمل ہوندا
اے۔ مناسب اخذی
شرطاں ہیٹھ ایسا
نظام اپنی بے
تضادی بیان کرن
والا جملہ آپ
ثابت نہیں کر
سکدا۔ اوس جملے
دی نفی وی ثابت
نہ ہو سکّن دا
دعویٰ وادھو
مفروضیاں اُتے
منحصر اے۔

ایہ نتیجہ ہلبرٹ دے
اصل منصوبے لئی
سخت ضرب سی۔ پر
جیویں ریاضی وچ
اکثر ہوندا اے،
ایہدے نال دلچسپ
نویں تحقیق دے
راہ وی کھُل گئے۔
جے ساری ریاضی
نوں سمیٹن والا
اکّو رسمی نظام
نہیں، تاں محدود
دائرے والے نظام
بنانے تے اوہناں
وچ کیہ ثابت ہو
سکدا اے، ایہ
پرکھّن دا کم
معنی رکھدا اے۔
ایہ وی فائدے دی گل اے
کہ انہاں نظاماں دی بے
تضادی ثابت کرن لئی
گھٹ پابندی والے طریقے
بنائیے تے ناپیئے کہ
اوہناں دی بے تضادی
ثابت کرنی کِنّی اوکھی اے۔
گوڈل نے دکھایا
کہ (تقریباً)
ہر رسمی نظام
کجھ سوالاں دا
فیصلہ نہیں
کردا؛ ایس
لئی کسے خاص
نظام دے
ایسے «دلچسپ»
غیر فیصلہ شدہ
سوال لبھنا،
تے ایہ جاننا
کہ اوہناں
دا فیصلہ کرن
لئی نظام
کِنّا طاقتور
ہونا لوڑیں دا
اے، اہم
اے۔ اج تک
منطق دان
ایہ سوال
ریاضی دی
اک نویں شاخ،
ثبوتاں دے
نظریے، وچ
کھوج رہے نیں۔

\end{document}
